% =====================================================================
\chapter{Distribuciones muestrales y teorema central del límite}
\label{tema:3}
% =====================================================================

Un estadístico es una variable aleatoria y, para hacer inferencia, hace falta conocer (al menos aproximadamente) su distribución. Este tema estudia la distribución de la media muestral: exacta cuando los datos son normales y aproximada, gracias al teorema central del límite, en cualquier otro caso. Se incluyen la aproximación normal a la binomial y las distribuciones $\chi^{2}$ y $t$ de Student, que se usarán en los intervalos de confianza y en los contrastes. Se sigue el capítulo~10 de Blitzstein y Hwang~\cite{bh} (§10.3 «Central limit theorem» y §10.4 «Chi-Square and Student-t»), la lectura~6b del curso 18.05 del MIT~\cite{mit1805}, «Central Limit Theorem and the Law of Large Numbers», y el capítulo~7 de OpenStax~\cite{openstax} (§7.1 y §7.2).

\section{La distribución normal}

\begin{definicion}[Distribución normal]\label{def:t3-normal}
Una variable $X$ tiene distribución \emph{normal} de media $\mu\in\R$ y varianza $\sigma^{2}>0$, y se escribe $X\sim\Nor(\mu,\sigma^{2})$, si su densidad es
\[
f(x)=\frac{1}{\sigma\sqrt{2\pi}}\exp\Bigl(-\frac{(x-\mu)^{2}}{2\sigma^{2}}\Bigr),\qquad x\in\R.
\]
La \emph{normal estándar} es $Z\sim\Nor(0,1)$; su función de distribución se denota $\Phi(z)=\PP(Z\le z)$ y su cuantil de orden $q$, $z_q=\Phi^{-1}(q)$.
\end{definicion}

La función $\Phi$ no tiene expresión mediante funciones elementales y se tabula (\cref{tab:phi}). Por la simetría de la densidad respecto de~$0$, $\Phi(-z)=1-\Phi(z)$ y $z_{1-q}=-z_q$.

\begin{proposicion}[Tipificación]\label{prop:t3-tipificacion}
Si $X\sim\Nor(\mu,\sigma^{2})$ y $a\neq0$, entonces $aX+b\sim\Nor(a\mu+b,a^{2}\sigma^{2})$. En particular, $Z=(X-\mu)/\sigma\sim\Nor(0,1)$ y
\[
\PP(X\le x)=\Phi\Bigl(\frac{x-\mu}{\sigma}\Bigr).
\]
\end{proposicion}

La demostración es el \cref{eje:t3-lineal-normal}. Con la tabla se obtiene la conocida regla empírica: una normal toma valores a menos de $1$, $2$ y $3$ desviaciones típicas de su media con probabilidades $2\Phi(1)-1\approx0{,}683$, $2\Phi(2)-1\approx0{,}954$ y $2\Phi(3)-1\approx0{,}997$.

\begin{teorema}[Combinaciones lineales de normales independientes]\label{teo:t3-suma-normales}
Si $X_1,\dots,X_n$ son independientes con $X_i\sim\Nor(\mu_i,\sigma_i^{2})$, entonces $\sum_ia_iX_i\sim\Nor\bigl(\sum_ia_i\mu_i,\sum_ia_i^{2}\sigma_i^{2}\bigr)$ para cualesquiera $a_i$ no todos nulos. En particular, si $X_1,\dots,X_n\sim\Nor(\mu,\sigma^{2})$ son \iid,
\[
\Xbar\sim\Nor\Bigl(\mu,\frac{\sigma^{2}}{n}\Bigr).
\]
\end{teorema}

La media y la varianza se siguen de la \cref{prop:t2-props}; que la distribución siga siendo normal se demuestra con funciones generadoras de momentos~\cite[§6.6]{bh} y se admite aquí sin demostración.

\section{El teorema central del límite}

Cuando los datos no son normales, la distribución exacta de $\Xbar$ suele ser complicada. El teorema central del límite afirma que, sin embargo, para $n$ grande es aproximadamente normal, sea cual sea la distribución de partida, con tal de que tenga varianza finita.

\begin{teorema}[Teorema central del límite]\label{teo:t3-tcl}
Sean $X_1,X_2,\dots$ \iid\ con media $\mu$ y varianza $0<\sigma^{2}<\infty$, y sea $S_n=X_1+\dots+X_n$. Entonces, para todo $z\in\R$,
\[
\lim_{n\to\infty}\PP\Bigl(\frac{\Xbar_n-\mu}{\sigma/\sqrt n}\le z\Bigr)=\lim_{n\to\infty}\PP\Bigl(\frac{S_n-n\mu}{\sigma\sqrt n}\le z\Bigr)=\Phi(z).
\]
\end{teorema}

Las dos variables tipificadas son la misma (basta multiplicar numerador y denominador por $n$). La demostración, con funciones generadoras de momentos, está en~\cite[§10.3]{bh}; aquí interesa su uso.

\begin{observacion}[Uso práctico]\label{obs:t3-uso}
Para $n$ «grande» se aproxima
\[
\Xbar\approx\Nor\Bigl(\mu,\frac{\sigma^{2}}{n}\Bigr),\qquad S_n\approx\Nor(n\mu,n\sigma^{2}).
\]
Una regla habitual es $n\ge30$~\cite[§7.1]{openstax}, pero la calidad de la aproximación depende de la forma de la distribución: con datos simétricos basta un $n$ pequeño y con datos muy asimétricos hace falta bastante más. El teorema exige varianza finita: para distribuciones de cola muy pesada, como la de Cauchy (o la $t$ con un grado de libertad), la media muestral no se estabiliza~\cite[§10.3]{bh}.
\end{observacion}

\begin{ejemplo}\label{ej:t3-tcl}
La longitud de las respuestas de un modelo tiene media $\mu=200$ tokens y desviación típica $\sigma=80$, con distribución asimétrica. Para una muestra de $n=64$ respuestas, el error estándar de la media es $80/\sqrt{64}=10$, y por el teorema central del límite
\[
\PP(\Xbar>215)\approx1-\Phi\Bigl(\frac{215-200}{10}\Bigr)=1-\Phi(1{,}5)=1-0{,}9332=0{,}0668.
\]
Una respuesta individual supera los $215$ tokens con mucha mayor probabilidad: la media de $64$ respuestas fluctúa ocho veces menos que cada respuesta.
\end{ejemplo}

\section{Aproximación normal a la binomial}

Si $X\sim\Bin(n,p)$, entonces $X$ tiene la misma distribución que la suma de $n$ variables $\Ber(p)$ \iid, con $\mu=p$ y $\sigma^{2}=p(1-p)$. El teorema central del límite da el resultado clásico de De Moivre y Laplace.

\begin{corolario}[Aproximación normal a la binomial]\label{cor:t3-binomial}
Si $X\sim\Bin(n,p)$ con $0<p<1$, para $n$ grande
\[
X\approx\Nor\bigl(np,\,np(1-p)\bigr),\qquad \widehat p=\frac Xn\approx\Nor\Bigl(p,\,\frac{p(1-p)}{n}\Bigr).
\]
Una regla habitual para usarla es $np\ge10$ y $n(1-p)\ge10$.
\end{corolario}

\begin{definicion}[Corrección por continuidad]\label{def:t3-cc}
Al aproximar una variable con valores enteros por una continua, cada valor entero $k$ se identifica con el intervalo $[k-\frac12,k+\frac12]$. Así, para enteros $a\le b$,
\[
\PP(a\le X\le b)\approx\Phi\Bigl(\frac{b+\frac12-np}{\sqrt{np(1-p)}}\Bigr)-\Phi\Bigl(\frac{a-\frac12-np}{\sqrt{np(1-p)}}\Bigr).
\]
\end{definicion}

\begin{ejemplo}\label{ej:t3-binomial}
Sea $X\sim\Bin(100;\,0{,}5)$: $np=50$ y $\sqrt{np(1-p)}=5$. Sin corrección, $\PP(X\ge60)\approx1-\Phi(2)=0{,}0228$; con corrección, $\PP(X\ge60)=\PP(X\ge59{,}5)\approx1-\Phi(1{,}9)=0{,}0287$. El valor exacto es $0{,}0284$: la corrección mejora mucho la aproximación.
\end{ejemplo}

\section{Distribuciones \texorpdfstring{$\chi^{2}$}{chi cuadrado} y \texorpdfstring{$t$}{t} de Student}

Cuando $\sigma$ es desconocida y se sustituye por $S$, la variable $(\Xbar-\mu)/(S/\sqrt n)$ ya no es normal. Su distribución exacta, para datos normales, se describe con dos familias derivadas de la normal~\cite[§10.4]{bh}.

\begin{definicion}[Distribuciones $\chi^{2}$ y $t$]\label{def:t3-chi-t}
Sean $Z_1,\dots,Z_k\sim\Nor(0,1)$ \iid\ La distribución de $V=Z_1^{2}+\dots+Z_k^{2}$ se llama \emph{$\chi^{2}$ con $k$ grados de libertad} y se escribe $V\sim\chi^{2}_k$. Si $Z\sim\Nor(0,1)$ y $V\sim\chi^{2}_k$ son independientes, la distribución de
\[
T=\frac{Z}{\sqrt{V/k}}
\]
se llama \emph{$t$ de Student con $k$ grados de libertad} y se escribe $T\sim t_k$. Sus cuantiles se denotan $\chi^{2}_{k;q}$ y $t_{k;q}$ (\cref{tab:t,tab:chi2}).
\end{definicion}

La $\chi^{2}_k$ toma valores positivos, tiene media $k$ y varianza $2k$ (\cref{eje:t3-chi-momentos}). La $t_k$ es simétrica respecto de~$0$ (\cref{eje:t3-t-simetrica}), tiene colas más pesadas que la normal y converge a ella cuando $k\to\infty$: por ejemplo, $t_{5;0{,}975}=2{,}571$, $t_{30;0{,}975}=2{,}042$ y $z_{0{,}975}=1{,}960$.

\begin{teorema}[Teorema de Fisher]\label{teo:t3-fisher}
Sean $X_1,\dots,X_n\sim\Nor(\mu,\sigma^{2})$ \iid\ con $n\ge2$. Entonces $\Xbar$ y $S^{2}$ son independientes, $\Xbar\sim\Nor(\mu,\sigma^{2}/n)$ y
\[
\frac{(n-1)S^{2}}{\sigma^{2}}\sim\chi^{2}_{n-1}.
\]
\end{teorema}

Se admite sin demostración~\cite[§10.4]{bh}. Los $n-1$ grados de libertad son los de la \cref{obs:t1-n-1}: la suma de cuadrados respecto de $\Xbar$ equivale a una suma de $n-1$ cuadrados de normales estándar independientes. Como $\E[\chi^{2}_{n-1}]=n-1$, el teorema reproduce, para datos normales, el resultado $\E[S^{2}]=\sigma^{2}$ del \cref{teo:t2-s2}, y como $\Var(\chi^{2}_{n-1})=2(n-1)$, da $\Var(S^{2})=2\sigma^{4}/(n-1)$, que se usó en el \cref{eje:t2-ecm-normal}.

\begin{corolario}[Estadístico $t$]\label{cor:t3-t}
En las condiciones del \cref{teo:t3-fisher},
\[
T=\frac{\Xbar-\mu}{S/\sqrt n}\sim t_{n-1}.
\]
\end{corolario}

\begin{proof}
Sean $Z=(\Xbar-\mu)/(\sigma/\sqrt n)\sim\Nor(0,1)$ y $V=(n-1)S^{2}/\sigma^{2}\sim\chi^{2}_{n-1}$, independientes por el \cref{teo:t3-fisher}. Entonces
\[
\frac{Z}{\sqrt{V/(n-1)}}=\frac{(\Xbar-\mu)\sqrt n/\sigma}{S/\sigma}=\frac{\Xbar-\mu}{S/\sqrt n},
\]
que tiene distribución $t_{n-1}$ por la \cref{def:t3-chi-t}.
\end{proof}

\begin{observacion}[Errores frecuentes]\label{obs:t3-errores}
(a) El teorema central del límite no dice que los \emph{datos} se vuelvan normales al aumentar $n$, sino que lo hace la distribución de su \emph{media} (o de su suma).
(b) Usar la desviación típica de la población, $\sigma$, en lugar del error estándar, $\sigma/\sqrt n$, al tipificar una media.
(c) Olvidar la corrección por continuidad con $n$ moderado o confundir su sentido: $\PP(X>59)=\PP(X\ge60)\approx\PP(Y\ge59{,}5)$, no $\PP(Y\ge60{,}5)$.
(d) Usar la $t$ para datos claramente no normales y $n$ pequeño: el \cref{cor:t3-t} exige normalidad; para $n$ grande, la $t$ y la normal prácticamente coinciden.
\end{observacion}

\begin{observacion}[Conexión con el curso]
El módulo~5 enuncia el teorema central del límite y lo aplica al método de Monte Carlo. En el módulo~7, el gradiente calculado con un minilote de $B$ ejemplos es una media muestral: su ruido decrece como $1/\sqrt B$. En los módulos~22 y~23, el error estándar de una exactitud medida sobre $n$ preguntas es $\sqrt{p(1-p)/n}$, que con $n=1000$ y $p=0{,}8$ vale unos $1{,}3$ puntos porcentuales; las diferencias de un punto entre modelos en un benchmark de ese tamaño están, por tanto, dentro del ruido de muestreo.
\end{observacion}

\begin{lecturas}
\item Blitzstein y Hwang~\cite{bh}, §5.4 «Normal», §10.3 «Central limit theorem» y §10.4 «Chi-Square and Student-t».
\item Orloff y Bloom~\cite{mit1805}, lecturas~6b «Central Limit Theorem and the Law of Large Numbers» y 5c «Gallery of Continuous Random Variables».
\item OpenStax~\cite{openstax}, cap.~7: §7.1 «The Central Limit Theorem for Sample Means (Averages)» y §7.2 «The Central Limit Theorem for Sums».
\item Wasserman~\cite{wasserman}, cap.~5 «Convergence of Random Variables».
\end{lecturas}

\section*{Ejercicios del tema 3}
\addcontentsline{toc}{section}{Ejercicios del tema 3}

\begin{ejercicio}{1}{eje:t3-vf}
Indique si cada afirmación es verdadera o falsa. (a) Por el teorema central del límite, si $n=1000$ los datos de una muestra de una exponencial son aproximadamente normales. (b) Si los datos son normales, $\Xbar$ es exactamente normal para todo $n$. (c) El error estándar de la media de $100$ observaciones es la décima parte de la desviación típica de cada observación. (d) La $t_{10}$ tiene más probabilidad en las colas que la $\Nor(0,1)$. (e) Si $X\sim\Bin(20;\,0{,}02)$, la aproximación normal es adecuada.
\end{ejercicio}

\begin{ejercicio}{1}{eje:t3-tabla}
Con la \cref{tab:phi,tab:zq}, calcule: (a) $\Phi(1{,}23)$; (b) $\PP(Z>-0{,}5)$; (c) $\PP(-1<Z<2)$; (d) $z_{0{,}95}$ y $z_{0{,}99}$; (e) el valor $c>0$ tal que $\PP(\abs Z\le c)=0{,}90$.
\end{ejercicio}

\begin{ejercicio}{2}{eje:t3-normal-basico}
Una puntuación $X$ sigue una distribución $\Nor(100,15^{2})$. Calcule $\PP(X>130)$, $\PP(85<X<115)$ y el percentil $90$ de $X$.
\end{ejercicio}

\begin{ejercicio}{2}{eje:t3-openstax-media}
Una población tiene media $90$ y desviación típica $15$. Se toma una muestra de tamaño $n=25$. (a) Calcule aproximadamente $\PP(85<\Xbar<92)$. (b) ¿Qué valor de $\Xbar$ está dos errores estándar por encima de la media?
\fuente{adaptado de OpenStax, \emph{Introductory Statistics}~2e~\cite{openstax}, §7.1, ejemplo~7.1.}
\end{ejercicio}

\begin{ejercicio}{2}{eje:t3-openstax-suma}
Con la población del \cref{eje:t3-openstax-media} y una muestra de tamaño $n=80$: (a) calcule aproximadamente la probabilidad de que la suma de los $80$ valores supere $7500$; (b) halle el valor de la suma que está $1{,}5$ desviaciones típicas por encima de su media.
\fuente{adaptado de OpenStax, \emph{Introductory Statistics}~2e~\cite{openstax}, §7.2, ejemplo~7.5.}
\end{ejercicio}

\begin{ejercicio}{2}{eje:t3-elecciones}
El $50\,\%$ de una población apoya al candidato A y el $30\,\%$ a otros candidatos menores. Se encuesta a $400$ personas elegidas al azar. Con el teorema central del límite, estime: (a) la probabilidad de que al menos el $52{,}5\,\%$ de los encuestados apoye a A; (b) la probabilidad de que menos del $31\,\%$ apoye a los candidatos menores.
\fuente{adaptado de Orloff y Bloom, MIT 18.05~\cite{mit1805}, \emph{problem set}~4, problema~2.}
\end{ejercicio}

\begin{ejercicio}{2}{eje:t3-redondeo}
Un comercio redondea el importe de cada pedido a los $5$ céntimos más próximos. Se supone que el error de redondeo de cada pedido toma los valores $-2,-1,0,1,2$ céntimos con igual probabilidad, de forma independiente. Con $1000$ pedidos diarios, estime la probabilidad de que el error total del día supere los $100$ céntimos en valor absoluto.
\fuente{adaptado de Orloff y Bloom, MIT 18.05~\cite{mit1805}, \emph{problem set}~4, problema~3.}
\end{ejercicio}

\begin{ejercicio}{2}{eje:t3-binomial}
Sea $X\sim\Bin(100;\,0{,}5)$ el número de caras en $100$ lanzamientos. (a) Aproxime $\PP(X=50)$ con la corrección por continuidad. (b) Aproxime $\PP(40\le X\le60)$ con y sin corrección. (c) Compare con los valores exactos $\PP(X=50)=0{,}0796$ y $\PP(40\le X\le60)=0{,}9648$.
\end{ejercicio}

\begin{ejercicio}{2}{eje:t3-benchmark}
Un modelo tiene exactitud real $p=0{,}8$ en un banco de preguntas muy grande. Se evalúa con $n=1000$ preguntas elegidas al azar. (a) Calcule el error estándar de $\widehat p$. (b) Estime $\PP(\widehat p\ge0{,}82)$. (c) Halle un intervalo centrado en $0{,}8$ que contenga a $\widehat p$ con probabilidad aproximada $0{,}95$.
\end{ejercicio}

\begin{ejercicio}{2}{eje:t3-tamano}
(a) ¿Qué tamaño muestral garantiza, según el teorema central del límite, que $\PP(\abs{\Xbar-\mu}<0{,}1\,\sigma)\ge0{,}95$? (b) Aplíquelo para estimar una exactitud con error menor que $0{,}05$ con probabilidad $0{,}95$, sea cual sea $p$. Compare con el \cref{eje:t2-chebyshev-n}.
\end{ejercicio}

\begin{ejercicio}{2}{eje:t3-t-cuantiles}
Con la \cref{tab:t}: (a) obtenga $t_{9;0{,}975}$, $t_{29;0{,}975}$ y $t_{120;0{,}975}$ y compárelos con $z_{0{,}975}$; (b) si $T\sim t_{15}$, calcule $\PP(T>2{,}131)$ y $\PP(\abs T\le1{,}753)$; (c) si $V\sim\chi^{2}_{10}$, halle $c$ con $\PP(V>c)=0{,}05$.
\end{ejercicio}

\begin{ejercicio}{3}{eje:t3-lineal-normal}
Demuestre la \cref{prop:t3-tipificacion}: si $X\sim\Nor(\mu,\sigma^{2})$ y $a\ne0$, entonces $Y=aX+b\sim\Nor(a\mu+b,a^{2}\sigma^{2})$. (\emph{Indicación:} calcule $\PP(Y\le y)$ distinguiendo $a>0$ y $a<0$, y derive.)
\end{ejercicio}

\begin{ejercicio}{3}{eje:t3-chi1}
Sea $Z\sim\Nor(0,1)$ y $V=Z^{2}$. Demuestre que $\PP(V\le v)=2\Phi(\sqrt v)-1$ para $v>0$ y deduzca que la densidad de la $\chi^{2}_1$ es
\[
f_V(v)=\frac{1}{\sqrt{2\pi v}}\,e^{-v/2},\qquad v>0.
\]
\end{ejercicio}

\begin{ejercicio}{3}{eje:t3-chi-momentos}
(a) Demuestre, integrando por partes, que si $Z\sim\Nor(0,1)$ entonces $\E[Z^{4}]=3\E[Z^{2}]=3$. (b) Deduzca que si $V\sim\chi^{2}_k$, entonces $\E[V]=k$ y $\Var(V)=2k$.
\end{ejercicio}

\begin{ejercicio}{3}{eje:t3-t-simetrica}
Sea $T\sim t_k$. Demuestre que $T$ y $-T$ tienen la misma distribución y deduzca que $t_{k;1-q}=-t_{k;q}$ y que $\PP(\abs T\le t_{k;1-\alpha/2})=1-\alpha$. (Se admite que la distribución de $T$ es continua.)
\end{ejercicio}
